\scanpage{50}
$\mathcal O_{\unsure{\underline\Sigma|\Sigma}}(1)$ sur $\Sigma$, défini par la représentation (10) de $\Sigma$ comme \deleted{fibré} \unsure{espace} projectif associé à un vectoriel bien déterminé.

Je désigne par $P\Sigma_{J,\mathbb C}$ ou simplement $P\Sigma$ le fibré principal associé de groupe $\mathbb C^*$ (en des espaces analytiques complexes, disons, avec sa topologie habituelle --- pour le moment on ne travaille pas dans le contexte des schémas\ldots). \deleted{\unsure{I}}

Un point $x$ de $\Sigma$ étant une droite $F$ de $V$ \deleted{\unsure{de $D$}}, on a des isom. can.
\[
(18)\qquad
\left\{\begin{aligned}
\unsure{\check V|\mathbb V}(\underline{\mathcal O}(1))_x&=\deleted{\unsure{\Delta}}V/F\\
F\otimes\mathcal O_F(1)&\simeq\det(V)\simeq\mathbb R(\underline\omega)\left(\overset{\mathrm{déf}}=\mathbb R\mathbin{\wedge}_{\{\pm1\}}\underline\omega\right)\\
\mathcal O_F(1)&\simeq\check F(\omega)\left(\overset{\mathrm{déf}}=\check F\mathbin{\wedge}_{\{\pm1\}}\underline\omega\right)\simeq F^\perp
\end{aligned}\right.
\]
\inserted{$\mathcal O_F(1)\overset{\mathrm{déf}}=$}
\margin{d'où pass. à dét}
\margin{\unsure{i.e.}}
\deleted{\unsure{$\mathbb R\otimes$}}
\margin{orth. pour la str. quadratique \unsure{invariante\ldots}}

d'où
\[
(19)\qquad(P\Sigma)_x=(P\Sigma)_F\simeq(V/F)^*\simeq(\check F)^*\mathbin{\wedge}_{\{\pm1\}}\underline\omega\simeq F^\perp
\]
\deleted{et par un} isom.\unsure{$^s$} compatibles avec l'action de $\mathbb C^*$ (opérant sur $\check F$ par contragrédience de son action sur $F$), on aura
\[
\deleted{\unsure{(19)}}\quad(20)\qquad(P\Sigma)_x\simeq F^*\mathbin{\wedge}_{\{\pm1\}}\underline\omega=(F\mathbin{\wedge}_{\{\pm1\}}\underline\omega)^*\simeq(F^\perp)^*
\]
\mspage{49}
